\chapter*{Introducción}
\addcontentsline{toc}{chapter}{Introducción}

El ajuste de controladores PID continúa siendo uno de los problemas más relevantes de la
ingeniería de control aplicada. Su vigencia no proviene de la complejidad de su estructura,
sino precisamente de lo contrario: el PID sigue utilizándose porque ofrece una relación
muy favorable entre simplicidad, interpretabilidad e impacto práctico. Sin embargo, esa
utilidad depende directamente de la manera en que se eligen sus ganancias, especialmente
cuando el sistema controlado exhibe no linealidades, restricciones físicas o perturbaciones
severas.

En microredes aisladas de base hidroeléctrica, como las que operan en distintas zonas
altoandinas del Perú, el problema de sintonización adquiere una dimensión crítica. La
dinámica de la turbina hidráulica, el efecto de golpe de ariete, la variación de la carga y la
intermitencia de la generación solar pueden degradar rápidamente el comportamiento de un
controlador sintonizado con reglas clásicas. Lo que en un banco académico luce aceptable,
en una planta real puede traducirse en oscilaciones de frecuencia, desgaste del actuador y
mayor vulnerabilidad del sistema ante eventos operativos cotidianos.

Frente a ese escenario, el presente libro aborda la sintonización PID desde una perspectiva
de optimización. La idea central es considerar que la elección de $K_p$, $K_i$ y $K_d$ no debe
depender solo de fórmulas cerradas o ajustes de prueba y error, sino de un procedimiento de
búsqueda orientado por métricas de desempeño claras. En este marco, los algoritmos
metaheurísticos ofrecen una alternativa atractiva porque permiten explorar espacios de
solución complejos sin exigir un modelo idealizado del problema.

La obra toma como hilo conductor el caso de una microred hidro-solar aislada en Limbani y
organiza la discusión en torno a cuatro niveles. Primero, se revisan los fundamentos del
control PID y los criterios con que se evalúa su desempeño. Segundo, se expone la forma de
plantear la sintonización como un problema de optimización restringida. Tercero, se estudian
varios métodos metaheurísticos con énfasis en PSO y su aplicación comparativa frente a
otros algoritmos de búsqueda. Finalmente, se discuten resultados, alcances y líneas futuras
de validación.

El recorrido del libro es el siguiente:
\begin{enumerate}[label=\arabic*.]
    \item los capítulos iniciales presentan los fundamentos del PID, los criterios de desempeño y las limitaciones de la sintonización clásica;
    \item los capítulos intermedios formulan el problema de optimización y desarrollan la base conceptual de los métodos metaheurísticos;
    \item los capítulos dedicados a PSO, GWO, Eagle Strategy y ABC describen su adaptación al problema de sintonización;
    \item los capítulos finales consolidan la metodología comparativa, los casos de estudio, la discusión de resultados y las conclusiones.
\end{enumerate}

Más que cerrar el tema, este libro busca dejar una plataforma de trabajo técnicamente
coherente, visualmente clara y lista para seguir creciendo con nuevos experimentos,
figuras, tablas y validaciones.
